% submission_info.tex (generated; edit src/submission_info.template.tex) -- what to enter on the submission server of the American Physical Society for the Regular Article
% "Higgs as a Composite, and Quantum Gravity" (Physical Review D).  PROJECT spacetimes, folder
% Revision_HS_CJT_and_Kohn-Sham/paper_prd.  Generated from src/submission_info.template.tex by scripts/build_paper.py,
% which inserts the abstract, the page counts and the table of the files from the build of the manuscript.
\documentclass[11pt]{article}
\usepackage[margin=2.5cm]{geometry}
\usepackage{amsmath,amssymb}
\usepackage{longtable}
\usepackage{array}
\usepackage{listings}
\usepackage[colorlinks=true,linkcolor=blue,urlcolor=blue]{hyperref}
\hypersetup{pdftitle={Submission information: Higgs as a Composite, and Quantum Gravity},pdfauthor={Patrick Lee Nash}}
\lstset{basicstyle=\ttfamily\footnotesize,breaklines=true,columns=fullflexible,frame=single,keepspaces=true,xleftmargin=4pt,xrightmargin=4pt}
\setlength{\parindent}{0pt}
\setlength{\parskip}{6pt}
\newcommand{\LL}{{\rm LL}}
\newcommand{\SF}{{\rm SF}}

\begin{document}

\begin{center}
{\Large\bf Submission information}\\[4pt]
{\large ``Higgs as a Composite, and Quantum Gravity''}\\[2pt]
Physical Review D, Regular Article\\[2pt]
PROJECT spacetimes, November 21, 2026
\end{center}

\section{What this document contains}

Everything that the submission server of the American Physical Society (APS) asks for, prepared from the manuscript: the
journal and the article type, the title, the byline, the abstract, the subject classification, the data availability
statement, the editorial information with the suggested referees, the cover letter, and the files to upload. The author
performs every step on the server himself; Sec.~\ref{sec:steps} lists the steps in order, and Sec.~\ref{sec:confirm} lists
the statements that only the author can confirm.

\section{Journal, article type and length}

\begin{itemize}
\item Journal: Physical Review D. Article type: Regular Article (not a Letter, Comment or Reply; no earlier paper needs to be
named).
\item Length: the manuscript has 10 pages in the two-column format of REVTeX 4.2, the Supplemental Material
36 pages, the cover letter 2 page(s). Physical Review D sets no limit to the length of Regular Articles
(the appropriate length depends on the information presented); the word limits of Letters and Comments do not apply.
\item Abstract: 335 words (a formula counted as one word); the APS guidance asks for less than 500 words.
\end{itemize}

\section{Title and byline}
\label{sec:title}

\begin{itemize}
\item Title: Higgs as a Composite, and Quantum Gravity
\item Author: Patrick Lee Nash, Ph.D. (the manuscript prints the name as the author asked; published Physical Review
bylines show names without academic degrees, so the production editors may print ``Patrick Lee Nash'').
\item Affiliation: PROJECT spacetimes. APS asks for the institution where the research was done, usually with the city,
the state and the country; the author completes it on the server (and in \texttt{\textbackslash affiliation} of the
source, if he wishes) with the address he wants printed.
\item Co-authors: none. ORCID: the server asks the author to link his ORCID iD.
\end{itemize}

\section{Abstract}
\label{sec:abstract}

As it appears in the manuscript:

\begin{quote}
The unit bi-octonions form a Moufang loop, the Lorentz super-loop, that contains $SL(2,\mathbb C)$ and eight spinor charges; these are carried by an isodoublet of DiracLike spinors on Minkowski space (Dirac spinors of $SL(2,\mathbb C)$ that have electric charge zero and no electroweak quantum numbers, particles not yet observed), coupled to a new gauge group $SL(2,\mathbb C)\times SU(2)_{\rm v}\times U(1)^{\rm F}$ of the local Lorentz transformations, a vector-like isospin and the fermion number. With a contact self-interaction whose spin-zero pair channel is supercritical, the spinor field condenses into a superfluid whose order parameter is a composite isospin vector of fermion number two. We solve it in five mean-field methods (Hubbard-Stratonovich with three decouplings, Kohn-Sham, and Cornwall-Jackiw-Tomboulis) and derive its effective Lagrangians. In every method the condensate is the Higgs field of the new group: the isospin bosons $W^\pm$ and the fermion-number boson $Z$ become massive, and exactly one gauge boson stays massless. This boson, which we name the zeugon, is not the photon: it couples to the unbroken isospin charge, it is the only gauge boson that attracts the two members of a pair, and, because the surviving group is $SL(2,\mathbb C)\times O(2)$, its charge conjugation is itself a gauge symmetry, the structure of Alice electrodynamics. The amplitude mode is a composite scalar bound below the threshold of two quasi-particles; once it can decay into the longitudinal gauge bosons its width is $2.1\mbox{--}3.9$ times its mass. From the static pair response we compute the gradient coefficient of the order parameter at the critical temperature and find Ginzburg numbers 7.6\mbox{--}11: the critical region is wider than the critical temperature itself, and the mean-field transition is not reliable. The Higgs particle of the new group is therefore composite and broad. Because the condensate is a singlet of the electroweak group, it cannot be the Higgs boson observed at the LHC. We outline how gravitation enters through E.~A.~Lord's gauge theory of the local change of spinor representation, for which the $SL(2,\mathbb C)$ factor of the gauge group is the natural place; the gravitational sector is not computed here.
\end{quote}

The LaTeX source of the same text, to paste into the abstract field of the form:

\begin{lstlisting}
The unit bi-octonions form a Moufang loop, the Lorentz super-loop, that contains $SL(2,\mathbb C)$ and eight spinor charges; these are carried by an isodoublet of DiracLike spinors on Minkowski space (Dirac spinors of $SL(2,\mathbb C)$ that have electric charge zero and no electroweak quantum numbers, particles not yet observed), coupled to a new gauge group $SL(2,\mathbb C)\times SU(2)_{\rm v}\times U(1)^{\rm F}$ of the local Lorentz transformations, a vector-like isospin and the fermion number. With a contact self-interaction whose spin-zero pair channel is supercritical, the spinor field condenses into a superfluid whose order parameter is a composite isospin vector of fermion number two. We solve it in five mean-field methods (Hubbard-Stratonovich with three decouplings, Kohn-Sham, and Cornwall-Jackiw-Tomboulis) and derive its effective Lagrangians. In every method the condensate is the Higgs field of the new group: the isospin bosons $W^\pm$ and the fermion-number boson $Z$ become massive, and exactly one gauge boson stays massless. This boson, which we name the zeugon, is not the photon: it couples to the unbroken isospin charge, it is the only gauge boson that attracts the two members of a pair, and, because the surviving group is $SL(2,\mathbb C)\times O(2)$, its charge conjugation is itself a gauge symmetry, the structure of Alice electrodynamics. The amplitude mode is a composite scalar bound below the threshold of two quasi-particles; once it can decay into the longitudinal gauge bosons its width is $2.1\mbox{--}3.9$ times its mass. From the static pair response we compute the gradient coefficient of the order parameter at the critical temperature and find Ginzburg numbers 7.6\mbox{--}11: the critical region is wider than the critical temperature itself, and the mean-field transition is not reliable. The Higgs particle of the new group is therefore composite and broad. Because the condensate is a singlet of the electroweak group, it cannot be the Higgs boson observed at the LHC. We outline how gravitation enters through E.~A.~Lord's gauge theory of the local change of spinor representation, for which the $SL(2,\mathbb C)$ factor of the gauge group is the natural place; the gravitational sector is not computed here.
\end{lstlisting}

\section{Subject classification (PhySH)}
\label{sec:physh}

The Physics Subject Headings are chosen by searching in the classification field of the form. The following concepts
were found in PhySH under \emph{Research Areas} $\to$ \emph{Hypothetical particle physics models} (discipline Particles
\& Fields): \emph{Composite models}, \emph{Dynamical symmetry breaking models}, \emph{Extensions of Higgs sector}. Suggested:
\begin{itemize}
\item primary concept (to be starred): \emph{Composite models};
\item further concept: \emph{Dynamical symmetry breaking models}. (\emph{Extensions of Higgs sector} is not suggested:
the Higgs field of the model is that of a new gauge group, whose spinor field has no electroweak quantum numbers, not an
extension of the Higgs sector of the Standard Model.)
\item further concepts to search for (their exact names in PhySH were not checked): dark or hidden sectors; new gauge
bosons; topological defects (Alice strings); spontaneous symmetry breaking; Ginzburg-Landau theory; critical phenomena;
mean field theory; finite-temperature field theory; superfluidity; quantum gravity.
\end{itemize}

\section{Data availability statement}
\label{sec:das}

Physical Review D publishes a data availability statement with every article, and the server asks for it with its
citations. The statement of the manuscript, with the citation written out:

\begin{quote}
The data that support the findings of this article are openly available in the repository of PROJECT spacetimes,
\url{https://github.com/once-ere/spacetimes}, folder \texttt{Revision\_HS\_CJT\_and\_Kohn-Sham}: the data files
(\texttt{data/}), the Rust and Python programs that produce them (\texttt{rust/}, \texttt{python/}), the reports of their
checks, the Wolfram Language notebook, and the scripts that build this article from the data (\texttt{paper\_prd/}). The
Supplemental Material lists the files with their SHA-256 checksums.
\end{quote}

A repository can change after publication. A frozen copy with a persistent identifier makes the statement permanent: for
example a release of the repository archived by a service that assigns a DOI (Zenodo archives releases of GitHub
repositories), whose DOI is then added to the statement and to the reference of the repository.

\section{Editorial information}
\label{sec:editorial}

\paragraph{Comments for the editor} (a proposal):
\begin{quote}
The Supplemental Material (file \texttt{higgs\_composite\_qg\_SM.pdf}) contains the derivations, the complete tables and
every numerical check with its deviation and tolerance; the programs and data are in a public repository (see the data
availability statement). This is not a joint submission. The manuscript builds on an unpublished manuscript of the author,
dated November 16, 2026, which can be provided on request.
\end{quote}

\paragraph{Suggested referees} (by field; the manuscript cites work of each). The editors choose the referees; the form
asks for names, institutions and e-mail addresses, which the author takes from the institutional web pages of the persons
he selects, after checking that none of them has a conflict of interest with him (a collaboration, a common institution,
a personal relation):
\begin{itemize}
\item Christopher T. Hill: dynamical symmetry breaking, top-quark condensation and the compositeness condition (W.~A.~Bardeen,
C.~T.~Hill and M.~Lindner, Phys.\ Rev.\ D 41, 1647 (1990); C.~T.~Hill and E.~H.~Simmons, Phys.\ Rep.\ 381, 235 (2003));
\item Roberto Contino: the composite Higgs boson as a pseudo-Goldstone boson (K.~Agashe, R.~Contino and A.~Pomarol, Nucl.\
Phys.\ B 719, 165 (2005));
\item Kari Rummukainen: the thermal electroweak transition on the lattice (K.~Kajantie, M.~Laine, K.~Rummukainen and
M.~Shaposhnikov, Phys.\ Rev.\ Lett.\ 77, 2887 (1996); M.~D'Onofrio and K.~Rummukainen, Phys.\ Rev.\ D 93, 025003 (2016));
\item Mikko Laine: thermal field theory and the electroweak transition (Phys.\ Rev.\ Lett.\ 77, 2887 (1996));
\item Andrea Pelissetto: critical phenomena and the renormalization group (A.~Pelissetto and E.~Vicari, Phys.\ Rep.\ 368,
549 (2002)).
\end{itemize}

\paragraph{Joint submissions, excluded referees:} none.

\section{Cover letter}

The file \texttt{cover\_letter.pdf} (2 page(s)) gives the context of the manuscript (the earlier manuscript
on the Lorentz super-loop), the four main findings with their numbers, the reasons for Physical Review D, a statement on
the title and the gravitational sector (a program, not computed), the reproducibility of the numbers, the submission
history (a first submission, not under consideration elsewhere) and a pointer to the suggested referees.

\section{Files to upload}
\label{sec:files}

The folder \texttt{paper\_prd/submission/} holds exactly these files. Its LaTeX source refers to the figures without a
folder, because the server keeps the uploaded files in one folder; the build compiles a copy of the folder alone and checks
that it gives the same number of pages without an undefined reference. The server builds its own PDF from the source,
which the author compares with \texttt{higgs\_composite\_qg.pdf} before submitting.

\begingroup\small\setlength{\tabcolsep}{4pt}
\begin{longtable}{>{\raggedright\arraybackslash}p{0.31\textwidth}>{\raggedright\arraybackslash}p{0.31\textwidth}r>{\raggedright\arraybackslash}p{0.2\textwidth}}
\hline\hline file & content & kB & SHA-256 (first 16 digits) \\ \hline
\endfirsthead
\hline file & content & kB & SHA-256 (first 16 digits) \\ \hline
\endhead
\hline\hline
\endlastfoot
\texttt{higgs\_composite\_qg.pdf} & the manuscript (PDF) & 561 & \texttt{f8f78ce49b38e0ed} \\
\texttt{higgs\_composite\_qg.tex} & the manuscript (LaTeX source, REVTeX 4.2) & 48 & \texttt{be51e0ee904836cb} \\
\texttt{higgs\_composite\_qg.bbl} & the bibliography of the source (from BibTeX with apsrev4-2.bst) & 29 & \texttt{05cf625687246687} \\
\texttt{fig\_part3\_q1\_potential.pdf} & Figure 1 & 78 & \texttt{d76e0840c72af9ee} \\
\texttt{fig\_part4\_r1\_massmatrix.pdf} & Figure 2 & 26 & \texttt{caa47847ba676acb} \\
\texttt{fig\_part3\_q4\_thermal.pdf} & Figure 3 & 43 & \texttt{c7bac6536492b220} \\
\texttt{fig\_prd\_ginzburg.pdf} & Figure 4 & 26 & \texttt{1e79446ef1c9c36e} \\
\texttt{fig\_prd\_higgs.pdf} & Figure 5 & 23 & \texttt{aaed0cb24621a722} \\
\texttt{higgs\_composite\_qg\_SM.pdf} & the Supplemental Material (PDF) & 635 & \texttt{68580df864e4b763} \\
\texttt{cover\_letter.pdf} & the cover letter (PDF) & 86 & \texttt{7a7af727c5d82a82} \\
\end{longtable}
\endgroup

\section{The steps on the submission server}
\label{sec:steps}

\begin{enumerate}
\item Log in to the APS submission server and choose ``Create a New Submission''.
\item Check or update the contact information at the top of the first page.
\item Select the journal Physical Review D and confirm the journal-specific guidelines.
\item Select the article type Regular Article.
\item Enter the title and the abstract (Secs.~\ref{sec:title} and \ref{sec:abstract}).
\item Co-authors: none; link the ORCID iD.
\item Subject classification: search the PhySH concepts of Sec.~\ref{sec:physh} and star the primary concept.
\item Open Science information: enter the data availability statement of Sec.~\ref{sec:das}.
\item Editorial information: the comments for the editor and the suggested referees of Sec.~\ref{sec:editorial}.
\item Upload the cover letter and the files of Sec.~\ref{sec:files} (the Supplemental Material as Supplemental Material).
\item Compare the PDF built by the server with \texttt{higgs\_composite\_qg.pdf}, then submit.
\end{enumerate}

\section{Statements that only the author can confirm}
\label{sec:confirm}

\begin{itemize}
\item The submission history of the cover letter: the manuscript has not been published or submitted elsewhere.
\item The status of the earlier manuscript on the Lorentz super-loop (unpublished, dated November 16, 2026), which the
cover letter and the comments for the editor state.
\item The absence of conflicts of interest with the suggested referees.
\item The affiliation and the address to be printed, and the form of the name.
\item The funding information that the server asks for (the manuscript names no funding).
\item The title: the manuscript and the cover letter say that the gravitational sector is a program and is not computed;
the editors or the referees may ask whether the title should say so.
\end{itemize}

\end{document}
